% Part: lambda-calculus
% Chapter: lambda-definability
% Section: lambda-definable-recursive

\documentclass[../../../include/open-logic-section]{subfiles}

\begin{document}

\olfileid{lam}{ldf}{ldr}
\olsection{توابع \usetoken{S}{lambda definable} بازگشتی‌اند}

نه‌تنها همهٔ توابع بازگشتی جزئی~!!{lambda definable} هستند، بلکه عکس
آن نیز درست است؛ یعنی همهٔ توابع~!!{lambda definable} بازگشتی جزئی‌اند.

\begin{thm}
  \ollabel{thm:lambda-computable} اگر تابع جزئی~$f$
  !!{lambda definable} باشد، بازگشتی جزئی است.
\end{thm}

\begin{proof}
  فقط طرح برهان را می‌آوریم. نخست ترم‌های~$\lambd$ را حسابی‌سازی
  می‌کنیم؛ یعنی با استفاده از کدگذاری متعارف دنباله‌ها به‌وسیلهٔ توان‌های
  اعداد اول، اعداد گودل را به‌طور نظام‌مند به ترم‌های~$\lambd$ نسبت
  می‌دهیم. سپس تابع بازگشتی جزئی~$\fn{normalize}(t)$ را تعریف می‌کنیم
  که عدد گودل~$t$ یک ترم لامبدا را به‌عنوان آرگومان می‌پذیرد و اگر ترم
  صورت نرمال داشته باشد، عدد گودل آن صورت را بازمی‌گرداند؛ وگرنه
  تعریف‌نشده است. سپس دو تابع بازگشتی جزئی~$\fn{toChurch}$ و
  $\fn{fromChurch}$ را تعریف می‌کنیم که اعداد طبیعی را به اعداد گودل
  عددنماهای چرچ متناظر و از آن‌ها نگاشت می‌کنند.

  با استفاده از این توابع بازگشتی، می‌توانیم~$f$ را به‌صورت تابعی
  بازگشتی جزئی تعریف کنیم. ترم~$\lambd$ای چون~$F$ وجود دارد که~$f$ را
  !!{lambda define}s. برای محاسبهٔ~$f(n_1, \dots, n_k)$، نخست با
  $\fn{toChurch}(n_i)$ اعداد گودل عددنماهای چرچ متناظر را به دست آورید؛
  سپس آن‌ها را به~$\Gn{F}$ بیفزایید تا عدد گودل ترم
  $F \num{n_1}\dots\num{n_k}$ حاصل شود. اکنون~$\fn{normalize}$ را بر
  این عدد گودل اعمال کنید. اگر~$f(n_1, \dots, n_k)$ تعریف‌شده باشد،
  $F \num{n_1}\dots\num{n_k}$ صورت نرمالی دارد که باید یک عددنمای چرچ
  باشد؛ وگرنه صورت نرمالی ندارد (و بنابراین
  \[\fn{normalize}(\Gn{F\num{n_1}\dots\num{n_k}})\]
  تعریف‌نشده است). سرانجام~$\fn{fromChurch}$ را بر عدد گودل ترم
  نرمال‌شده اعمال کنید.
\end{proof}

\end{document}
